% s-202-p202-m001; original PDF 202
\Aphorism{s-202-p202-m001}{287}
Wie selten erlebt man als Autor, von \Sperr{demselben} die Feder wegzulegen, zu dem man sie angesetzt hat! Es ist ein tiefes Missverhältniss zwischen dem Gedachten und Geschriebenen, zwischen Wurf und Ziel, Plan und Gelingen. Der Gedanke färbt sich um, indem man ihn festhält, wie einer jener fabelhaften Thierkörper, die beim Nahen von Freund oder Feind ein Farbenspiel sehen lassen: sträubt er sich, dieser Gedanke, oder freut er sich, gegriffen und geformt zu werden? Hier ist ein Geheimniss, und ein Reiz des Autorseins, — bisweilen auch eine Qual, zum Beispiel, wenn der Gedanke sich eigensinnig gegen die Gestaltung in neuen, noch unversuchten Worten wehrt und durchaus in schon befahrene Geleise einbiegen will. Zwischen Gehirn und Hand, Feder und Papier geht viel verloren, und oft das Eigenthümliche, noch nicht Dagewesene der Conception; immer ist viel Affinität des neuen Weins zum alten Schlauch zu überwinden.
\par

